% Guide paragraphs written incrementally after image collation of each section.
\paragraph{110. Schematization through Postnikov fibrations, and its obstruction (pp. 443--445; 134-7/1--3).}
Grothendieck resumes the question of homotopy types over a ground ring by comparing an abelian group $M$ with the functor $W(M)(k)=M\otimes_{\underline Z}k$: the discrete object acquires a structure analogous to a vector bundle, suggesting an analogous construction for semisimplicial models and successive Postnikov fibers. The proposed ``schematization above level n'' would leave a lower truncation discrete and require the comparison from schematic cohomology with quasi-coherent coefficients to ordinary cohomology to be bijective. The reference to section 94 situates the hope of replacing $\underline Z$ by $\underline Q$ when nilpotence replaces an abelian fundamental group. On p.~444 he immediately questions this construction: torsion and the universal coefficients formula prevent the homotopy groups after base change from being obtained by tensoring the individual homotopy groups. The alternative is to reconstruct a homotopy type from a suitably enriched ``abelianization'' or ``linearization''; section 111 then seeks a different construction without explicit Postnikov decomposition. The handwritten replacement of ``(= homotopy)'' by ``(or rather, homotopy)'' on p.~444 matters here: the arXiv witness retains the cancelled equality wording. Its mapsto in the displayed base-change formula also differs from the scan's ordinary arrow; the scan alternates star subscripts with the typed notation X. Page 445 closes with a methodological reflection on the tension between detailed preparatory work and the desire to move quickly. The rust-obscured words on p.~443 remain marked rather than restored from the mathematical argument.

\paragraph{111. Unipotent bundles, divided powers, and a revised notion of equivalence (pp. 446--459; 134-7/4--17).}
The starting problem is to construct categories $Hot_n(k)$ of schematized $n$-connected homotopy types, with extension and restriction of the ground ring and a comparison with the ordinary homotopy category over $\underline Z$ (pp.~446--449). Grothendieck builds $U(k)$ inductively from affine bundles $W(M)$ and torsors, first in the fpqc setting. He asks whether the successive fibration data belong to the object or merely witness its existence, then proposes semisimplicial objects subject to connectivity and filtration conditions. The proposed weak equivalences are initially detected by taking sections, and the hoped-for $k$-module structures on homotopy groups lead to a search for a linearization into $D.(Ab_k)$. On p.~448 the model-category symbol has a typed subscript $m$, although the surrounding development uses $n$; this apparent index slip is retained in the text.

Pages 450--455 test divided powers as the required linearization. The candidate universal map $x\mapsto\sum x^{(i)}$ immediately forces a completion, and the completed divided-power coalgebra is considered through its group-like elements. After the Eilenberg--Maclane test, the discussion returns to pointed objects and uncompleted divided powers. The comparison (w) on p.~454 places ordinary free abelian linearization, $L$, and the pointed version $L_{pt}$ in one diagram. The marginal correction on p.~452 is consequential: ``One has to take k=Z, and components $X_n$ projective.'' Both its initial words and the subscript are legible on the scan. The subsequent distinction between quasi-isomorphisms and chain homotopisms supplies the projectivity restriction. The final reference in the sentence on p.~455 is to (x''), after a preceding reference to (x'''); the arXiv witness repeats (x''') and loses this distinction.

The argument is substantially revised on pp.~456--459. The author retracts the unrestricted derived-functor claim of p.~454, requires projective or flat components, and remarks that the Postnikov approach must again be set aside. The plain typed $I_k(M_*')$ on p.~456 is retained as printed, although the neighboring formulas suggest a notation slip. Page 457 then exposes the problem with ordinary homology as a $k$-module: the exterior-power calculation for $K(M,1)$, applied to $k=\underline Z[T]/(T^2+1)$, yields an obstruction. Page 458 accordingly defines total linearized homology $LH.(X_*)=L_{pt}(X_*)$ and considers weak equivalences detected by its homology invariants. The handwritten revision says it is becoming \emph{dubious} whether equivalences should be defined through $X_*(k)$; the arXiv witness reads \emph{obvious}, reversing the author's change of judgment. At the end of the section, Question 2 asks whether the two definitions coincide for projective pointed unipotent bundles satisfying $X_0=X_1=e$ (p.~459). The weaker condition $X_0=e$ raises an additional issue about abelian fundamental groups. These pages formulate questions and revised expectations, rather than establish an unrestricted equivalence theorem; section 112 resumes from the projective $\underline Z$ case.

\paragraph{112. Kan conditions and the scope of the linearization program (pp. 460--465; 134-7/18--23).}
The entry dated 26.8 takes the projective $\underline Z$ case of the comparison maps (w) as its working foothold. Grothendieck sketches a passage from the Eilenberg--Maclane case to finite chain complexes, then through finite and ``way-out'' infinite successive fibrations, using multiplicativity and hoped-for spectral-sequence or derived-category transitivity statements (pp.~460--461). The proposed Kan--Postnikov complexes use shifted projective resolutions of homotopy groups. To transfer the construction from discrete to schematic objects, he returns explicitly to the cohomology obstruction of p.~444 and hopes to dualize the homology comparison. This is an outline requiring substantial verification, as the author himself stresses. Constant complexes then show why $X_0=e$ is necessary, and the monoid discussion turns the issue into comparison of polynomial with arbitrary cochains. The multiplicative-monoid example on p.~462 is offered as a counterexample to the assumptions previously used; the circled marginal cross is part of the page's evidence and is absent from the arXiv witness.

The response is to ``drop Postnikow and keep Kan'' (p.~462). Since $U(k)$ is not closed under arbitrary fiber products, a schematic horn condition must ensure both representability of the relevant finite limits and a suitably strong surjectivity property (p.~463). The tentative name ``submersions'' covers alternatives ranging from trivial bundles to formally smooth maps admitting sections. The first nontrivial condition asks for the three maps $X_2\to X_1\times X_1$ to be submersions; for a monoid nerve this is meant to force a group. Page 464 lists the desired consequences: the linearization theorem over $\underline Z$, a homology criterion for homotopisms of connected Kan complexes over a general ring, homotopy fibers and Postnikov decomposition, and eventually equivalence of $Hot_n(\underline Z)$ with ordinary $Hot_n$. Page 465 distinguishes restriction of rings from extension, noting that restriction does not preserve projective components without an additional assumption and suggesting resolutions to remove that restriction. Its last sentence proposes compatibility of ring extension with the individual homotopy invariants; the next entry immediately revises that proposal. This transition is essential to reading the section as a developing argument rather than a list of established properties.

\paragraph{113. Soft Postnikov decomposition and a schematic fundamental group (pp. 466--472; 134-7/24--30).}
The entry dated 27.8 corrects the previous evening's assertion. Formula (A) concerns total homology and derived extension of scalars; the analogous termwise assertion (B) for homotopy modules fails without flatness hypotheses (p.~466). The revised arrow in (B) points left, a direction separately emphasized by a circled marginal arrow omitted by the arXiv witness. The reason is the use of shifted projective resolutions in a Postnikov fiber: after tensoring, higher homology can appear, so the base-changed successive fibration need not remain a Postnikov decomposition (p.~467). This explicitly answers the puzzle of p.~444 and motivates the judgment that ordinary Postnikov decomposition is not ``geometric,'' meaning compatible with extension of the ground ring.

The proposed alternative keeps the entire abelianization complex. Using the completed divided-power linearization and the embedding $X\to W(L(X))$, Grothendieck proposes repeatedly taking the homotopy fiber of $X_*\to W(L(X_*))$, without first extracting the first nonzero homology module (pp.~467--468). He names this ``soft Postnikov d\'evissage.'' For a discrete $K(G,1)$ the analogy leads to iterated commutators, and he distinguishes the versions associated with solvability and nilpotence. Formula (D) contains an index problem: the visible inner exponent is $n+1$, like the outer exponent, while the arXiv witness gives $n-1$. The checked text preserves the scan's reinforced plus and does not silently substitute a mathematically expected recurrence. The next page adds the issue of the action of $\pi_1$ on higher homotopy groups; trivial or unipotent action appears relevant to unipotent schematization (p.~469).

The second part of the section revisits complexes with $X_0=e$ but possibly nontrivial $X_1$. A group object gives a fundamental group that need not be abelian, and finite-presentation and smoothness assumptions lead to questions about nilpotence (p.~469). The Kan relations then motivate the presheaf $k'\mapsto\pi_1(X_*(k'))$. The desired invariant is a group object of $U(k)$ representing this sheaf, with an underlying parametrization by a module and a polynomial group law (pp.~470--471). This changes what the notation $\pi_1$ denotes: the group of sections alone is no longer the complete object. Formulas (G) and (G') express evaluation and base-change compatibility for the proposed schematic group. The following duplicated equation on p.~471 is nevertheless literally tautological in the scan, with $\pi_1(X_*(k))$ on both sides; the witness edits one side. The blank between $k$ and $k'$ in the preceding prose likewise contains no arrow on the page.

To admit arbitrary quotient modules, the notion of submersion is reconsidered once more: maps isomorphic to those induced by epimorphisms of modules are proposed as the most convincing of the three possibilities considered so far (p.~471). The printed ``unless \ldots{} does'nt give rise'' contains an apparent extra negation, retained despite the witness's change to ``gives rise.'' The stronger requirement that the schematic fundamental group come from addition on a module is distinguished from commutativity; the discussion ends with tests involving the linearization theorem, a fibration over $K(G,1)$, and the relationship between nilpotent discrete groups and group objects of $U(k)$ (p.~472, referring back to section 94). The handwritten addition about unipotent action on the higher $\pi_i$ belongs here and is not merely incidental: it qualifies the contemplated range of homotopy types. Section 114 begins on the same page by describing the accumulated picture as provisional and preparing a review of the program.

\paragraph{114. A provisional model category and the completion gap (pp. 472--480; 134-7/30--38).}
This review arranges the preceding experiments into a prospective construction. The category $U(k)$ consists of unipotent bundles built from modules; smooth objects correspond to projective modules, submersions to epimorphisms, and smooth submersions to maps with projective kernel (pp.~472--473). The completed divided-power coalgebra describes the unpointed linearization, while the pointed version removes the completion in the basic construction (pp.~474--475). Completion is a substantive distinction: p.~474 explicitly says that the map (C) of p.~467 ``does'nt quite exist'' with $W$ as written, and introduces the completed version $W^{\wedge}$ instead. Base change for the completed functor consequently uses completed tensor products. The subsequent model categories impose $X_i=e$ for $i\leq n$, take sections to pointed semisimplicial sets, and localize at the proposed weak equivalences. The hoped-for comparison with ordinary homotopy types is restricted to the appropriate connected range over $\underline Z$ (p.~476). On p.~477 the first expression for the category of models visibly uses $\pi_n(k)$, although the surrounding formulas use the handwritten model-category symbol $\mathcal M_n(k)$; this notation anomaly is retained rather than silently regularized.

Smooth Kan resolutions are then proposed to derive extension of scalars and the total linearization $LH.(X_*)=L_{pt}(X_*)$ (p.~477). The next two pages ask whether this homology functor is conservative and whether pointed linearization and completed linearization are quasi-isomorphic. They also reconsider comparison with ordinary free-module homology. The second question undergoes an immediate correction: on p.~479 the author says the proposed comparison with $k^{(X_*(k))}$ is generally false away from $k=\underline Z$. In the first nonvanishing degree, the expected $k$-module invariant is identified with integral homology of the section complex; tensoring that homology with $k$ produces a different candidate. The word \emph{integral} is explicit in the scan, whereas the arXiv witness gives \emph{internal}. The correction renews the need to construct homotopy modules intrinsically and to obtain them through a Postnikov argument, rather than transfer all ordinary homology assertions indiscriminately.

The final page makes the outstanding construction precise. The exponential map gives $X_*\to W^{\wedge}(L(X_*))$, but the required homotopy-class map is $X_*\to W(L_{pt}(X_*))$ (p.~480). Even a quasi-isomorphism between the two linearizations does not by itself remove the difference between $W^{\wedge}$ and $W$; the handwritten phrase ``flaw in this `argument''' marks exactly this gap. The target of the completed map is a semisimplicial sheaf and need not belong to the proposed model category. Hard or soft Postnikov decomposition therefore still requires this factorization, a suitable notion of homotopy fiber, and homotopy constructions internal to the model category. The failure of $U(k)$ to have fiber products explains why Quillen's standard machinery cannot simply be applied as stated. These are explicit unresolved tasks, and section115 returns the following morning to the meaning of the coefficient sheaf itself.

\paragraph{115. Quasicoherent coefficients and prorepresentable linearization (pp. 481--489; 134-7/39--47).}
The entry dated28.8 restarts the formula numbering and identifies the source of the preceding confusion. The appropriate free linearization is over the quasicoherent ring sheaf $\underline O_k$ on the fpqc site, not over a constant ring (pp.~481--482). Grothendieck explains that the obstruction groups in schematic Postnikov decomposition use the quasicoherent sheaf associated with a homotopy module, rather than the same module regarded as a constant coefficient group. The circled marginal clarification of the argument $X(n)_*,\pi_{n+1}$ belongs to this distinction; the overwritten version in the running formula remains marked as doubtful. He connects the earlier mistake to having translated geometric language into commutative algebra too thoroughly, losing sight of which coefficient objects were intended.

The universal free Module $\underline O_k^{(X)}$ gives the Hom, Ext and derived-Hom expressions (4)--(5') (p.~483), but is generally not quasicoherent. The proposed replacement would represent the same functor on quasicoherent coefficients. For $X=W(M)$, the continuous-Hom calculation through truncated divided-power modules shows that the correct object is a proobject, rather than the completed module regarded as an ordinary module (pp.~484--485). The obstruction is again scalar extension: inverse limits do not generally commute with tensoring by a new ground ring. Formula(13) packages the finite products of divided powers into a strict projective system; ordinary representability would require it to become essentially constant. The added phrase ``but with slightly different meaning'' therefore changes the interpretation of $L$ in (14), now valued in pro-modules. The pointed functor still takes values in ordinary modules, embedded as constant proobjects. Formula(18), on p.~486, supplies the exponential map into $WL(X)$ in the procategory and distinguishes this abelianization from its richer linear structure.

The proposed derived comparison (20) and its Ext form are then tested on $X=W(M)$ (pp.~486--487). Formula(21) contains a meaningful handwritten revision from a limit expression to a direct sum of Ext groups of divided powers. Projectivity of $M$ gives the expected vanishing in positive degree, with a passage to finite-type submodules sketched for the infinite case. For nonprojective $M$, the potential nonzero Ext$^1$ term challenges the author's previous expectation that quasicoherent cohomology of these bundles should vanish. His reference to Larry Breen is a question about this expectation, not a proof of the general formula. The text visibly says ``quasiunipotent'' where the witness shortens it to ``unipotent'' (p.~487). The next page restricts attention to the smooth case and postpones the projective-resolution argument for full generality.

For a semisimplicial complex, (22) uses Hyperext rather than a degreewise Ext construction; (23) compares the resulting pro-linearization with pointed linearization (p.~488). The comparison map(24) revisits the map from p.~454 with its proobject interpretation made explicit. Its relationship to a quasi-isomorphism, and then to the componentwise inverse limit used on p.~478, still requires foundational work. Page489 plainly admits that no argument has yet been found for the desired assertion under the proposed smoothness, connectivity and possibly Kan hypotheses. The second comparison is reformulated as $\underline O_k^{(X_*)}\to L(X_*)$ in (25). The author does not confidently assert that this is a quasi-isomorphism: quasicoherence of the first homology sheaves and essential constancy of the second remain in doubt. The practical hope is instead that it gives the needed comparisons against quasicoherent coefficients through(20). Section116 subsequently reports the collapse of the earlier general abelianization theorem, so these provisional claims must be read with that later revision in view.

\paragraph{116. The failed abelianization theorem and the interruption of the notes (pp. 490--491; 134-7/48--49).}
The entry dated 7.9 reports the decisive revision of the preceding program. After finding a simple description of schematic homotopy groups, Grothendieck learned in a telephone conversation with Illusie that the asserted comparison between discrete and schematic homology or cohomology was false. The anticipated equivalence between schematic and ordinary homotopy types therefore fails (p.~490). This is an explicit retraction of a central assumption underlying sections111--115; it does not cause him to abandon schematic homotopy types, whose formal properties he now wants to study in their own right, possibly through different models. The remainder of the section explains the intervening pause in writing through fatigue and the need to attend to the death of his granddaughter Ella. Page491 distinguishes work allowed to take place through attention from a compulsive forward movement that leaves perceptions disconnected from one's picture of reality. The revisions matter to the voice of this passage: the hand moves ``who'' after ``finally,'' retains ``his language'' for the body's intervention, and changes ``cause'' to ``caused.'' The witness substitutes ``which'' and ``its'' and loses the last correction. These pages provide both the mathematical warning needed to read the preceding conjectures and the author's own account of why the notebook paused before the technical review resumes in section117.

\paragraph{117. Size, exactness, and the distinction between schemes and formal schemes (pp. 492--496; 134-7/50--54).}
The entry dated 10.9 begins with two technical qualifications. The affine site used earlier is not automatically small in the working universe; a possible remedy is to bound the sizes of the modules and the rings used as arguments, using the fppf site for finite-type bundles (p.~492). More consequentially, $M\mapsto W(M)$ is right exact but need not preserve kernels. Thus a free resolution of $M$ need not become a resolution of $W(M)$. The derived-Hom comparison used to interpret formula(21) of p.~487 is therefore questioned in general and retained for the flat case (p.~493). This is a specific dependency of the cohomology calculation, not merely a set-theoretic caveat.

The finite projective case clarifies the geometry through duality. Polynomial maps can be described by maps of symmetric algebras or by continuous maps of the dual completed coalgebras; preserving the origin translates into augmentation compatibility (pp.~493--494). The uncompleted divided-power coalgebra is instead dual to the adically completed symmetric algebra. Consequently, the two algebraic descriptions lead to schemes and formal completions along their zero sections, respectively. The scan's formulas require care: the prose on p.~493 includes a dual check on $M$, but formulas(2), (3) and (4) omit dual signs that the arXiv witness supplies. Those omissions remain in the checked text. The upward-arrow-shaped completion mark in(3') on p.~494 is likewise recorded as drawn, and a separate unreadable pen mark above the star in(5) remains marked.

This geometric interpretation overturns the earlier hoped-for equivalence between the pointed and nonpointed linearizations (pp.~495--496). For the additive group in characteristic $p$, the degree-one comparison distinguishes additive polynomials from formal series $\sum_{i\geq0}c_i t^{p^i}$: the schematic side allows only finitely many nonzero coefficients. Grothendieck acknowledges that this example is neither simply connected nor over $\underline Z$, but states that he also no longer expects the comparison to be a quasi-isomorphism for the basic connected examples previously contemplated. He reports that Illusie rejected the vanishing assumptions behind his attempted ``way-out'' argument. The closing underline therefore warns explicitly against substituting $L_{pt}(W(M))$ for $L(W(M))$ when computing schematic homology or cohomology. Formal invariants have their own meaning and are distinct. Section118 begins on the same page with a different and deliberately striking claim: homotopy invariants themselves should nevertheless be recoverable from the first-order tangent object.

\paragraph{118. The tangent complex as a proposed homotopy invariant (pp. 496--500; 134-7/54--58).}
Against the distinction just established between schematic and formal cohomology, Grothendieck proposes that homotopy groups should depend only on the first-order tangent object. The Lie functor sends a smooth pointed bundle $W(M)$ to $M$, and formula(9) identifies the homotopy groups of $X_*(k)$ with those of the semisimplicial module $\mathrm{Lie}(X_*)$, hence with homology of its associated chain complex (p.~496). The next page sharply qualifies the claim: he has no direct description of a natural isomorphism under the stated hypotheses. It is tautological for complexes arising from chain complexes and is expected to extend through Postnikov decomposition, if every relevant model can be replaced by one admitting such a decomposition; naturality would still require work (p.~497). The typed revision makes this formula a ``cornerstone'' of a theory of schematic homotopy types and a proposed intrinsic definition of the invariants, rather than subordinating them from the outset to discrete homotopy groups. Weak equivalences would then be detected by the tangent complexes. Both subsequent comparison displays map an $X_*$ object to a $Y_*$ object; the arXiv witness repeats $X_*$ as target in each.

The contemplated application is to explicit models of spheres (pp.~498--499). To make the Lie formula useful for ordinary homotopy groups, one must first realize the discrete type schematically over $\underline Z$ and obtain a workable model. Successive suspension from $S^1$ suggests the need for suspension constructions, complementing the earlier need for fibers and loops; the proposed derivator would organize the finite integration and cointegration operations. These are motivations and difficult construction problems, not completed calculations of sphere homotopy groups. The discussion of $S^2$ as a homogeneous space under $S^3$, and of the corresponding group/flag-variety geometry, asks for a model of the entire relationship: a group, a subgroup playing the role of a torus or Borel subgroup, and a quotient. Several overprinted group-name letters remain uncertain in the text, while the star on $\underline C^*$ is visibly an asterisk rather than the witness's multiplication cross.

Hemispherical complexes supply a second route (pp.~499--500). Their smaller list of boundary and degeneracy operations, together with a prospective direct computation of homotopy in the discrete setting, is compared with the tangent-complex method. The marginal exclamation and vertical rule emphasize this alternative. The author immediately admits that he has not checked the proposed computation for arbitrary hemispherical sets that do not arise from an infinity-groupoid or infinity-stack. Comparing the two descriptions is therefore offered as a possible route toward understanding formula(9), with that unresolved premise preserved. Section119 then resumes on26.9 after another interruption and returns to the failure of the earlier homology comparison.

\paragraph{119. Integer-valued polynomials and a different comparison problem (pp. 500--505; 134-7/58--63).}
The entry dated 26.9 returns to the failed comparison of discrete and schematic linearization. Grothendieck reports Illusie's objection from Breen's work on stable cohomology, while admitting that the precise counterexample has not yet been made explicit to him (p.~500). He then gives a concrete degree-two obstruction using central extensions of the additive group (p.~501). The multiplicative group of truncated series $1+at+bt^2$ defines the extension $W_2$; a polynomial splitting requires $P(a+a')=P(a)+P(a')+aa'$. The solutions displayed over $\underline Q$ involve $t(t-1)/2$, which has no integral-coefficient representative after addition of the displayed linear term. This separates polynomial cochains over $\underline Z$ from arbitrary cochains on the discrete group. The final splitting display literally has $t$, without the square appearing in the preceding general homomorphism display; the checked text preserves that visible slip rather than adopting the arXiv witness's $t^2$. The main argument should therefore be read through the preceding coefficient equation, not by treating the last display as an independently correct formula.

The next proposal replaces polynomial algebras by ``binomial algebras,'' built from the integer-valued expressions $P_n(t)=t(t-1)\cdots(t-n+1)/n!$ (p.~502). Grothendieck tentatively requires the ground ring to carry the corresponding binomial operations and compares this enrichment with divided powers in crystalline cohomology. The marginal reference to pp.~506--507 anticipates the next section's reconsideration of these structures and of the restrictions on the ground ring. On p.~503 the finite-type qualification on the free module $M$ is an above-line typed insertion. The algebra $\mathrm{Symbin}_{\underline Z}(M)$ is defined through rational polynomial functions integral-valued on the dual lattice, and the map to all functions on a semisimplicial free abelian group is posed as a possible quasi-isomorphism in the connected range. Reduction to $K(\underline Z,n)$ is proposed by d\'evissage, not proved. The broader question asks whether Eilenberg--Maclane homology or cohomology is derived from a suitable nonadditive functor $B$, leading to ``B-schematic homotopy types'' and unipotent $B$-bundles. The dark superscript on the first Symbin display remains unresolved and is not silently discarded.

The author declines to pursue that alternative immediately, because the ordinary schematic program still promises a simple Lie-functor description of homotopy invariants (pp.~503--504). Loops, suspensions, and compatibility with passage to discrete types remain the practical tests. Reading Bousfield--Kan then changes the comparison: their construction is identified as $k$-completion, distinct from a fully developed category of types over $k$. The marginal ``definitely not, in general!'' rejects the tentative identification of completion with extension of the ground ring; the faint underline beneath ``ring extension'' identifies the disputed phrase (p.~504). Page~505 links this reading to the integration and cointegration operations of section~69 and to Don Anderson's existence theorem in closed model categories. A pen transposition moves the closed-model-category qualification directly after ``operations,'' making the setting of that theorem clearer. The section closes by recognizing Bousfield--Kan as foundational reading while postponing it long enough to record the recent schematic reflections. Section~120 starts on 28.9 by abstracting the extra algebraic operations just proposed.

\paragraph{120. Extra algebraic operations and polyhedral models (pp. 506--507; 134-7/64--65).}
The entry dated 28.9 abstracts the previous section's binomial proposal together with divided powers, lambda operations, and Adams-type operations. Grothendieck uses the notion of an ``analyzer'' and asks for an extension of the analyzer of commutative rings in which unary operations determine higher components by tensor powers. The handwritten changes from $n$ to $n+1$ tensor factors and from the older bounds to $n\geq-1$ and $n\geq1$ affect the indexing of the proposal (p.~506). The live corrected bounds are clear, although overwritten earlier digits remain doubtful. The essential datum is described as a ground ring $k_0$ and an algebra $\Omega_0$ equipped with composition. An $\Omega$-structure on a ring is then an action of these operations satisfying compatibility with addition, multiplication, and constants. Page~507 displays the two diagonals that encode the first two conditions. This is an outline of the desired axioms; the author expressly declines to write all of them down. The circled instructions ``\`a la ligne'' and ``id'' on p.~506 are retained as typist instructions, whose precise targets are not sufficiently explicit to justify speculative paragraph restructuring.

The second half of p.~507 turns from these operations to David W.~Jones's work on polyhedral cells and poly-T-complexes, supplied by Ronnie Brown. Grothendieck relates it to the questions of sections~91 and~93: when does a category of cells form a weak test category, and when does the proposed standard chain complex compute the homology of its presheaves? He finds that Jones's emphasis is instead on generalizing thin structures and a corresponding equivalence of categories. The original questions therefore remain open for this notebook. Both appearances of the presheaf-category notation have a hand-drawn upward caret in the scan, rather than the downward vee supplied by the witness. Section~121 next returns to the formal role that linearization should play across the various possible schematic approaches.

\paragraph{121. Abelianization, adjunction, and the pointed correction (pp. 508--510; 134-7/66--68).}
Grothendieck begins the formal review by defining $\mathrm{Hotab}=D.(Ab)$ and considering abel\-ian\-iza\-tion $LH.$ together with the Eilenberg--Maclane functor $K_\pi$ (p.~508). He relates their adjunction to the free-abelian-group and forgetful functors on semisimplicial objects, then interprets ordinary cohomology through the representing objects $K(\pi,n)$. The notation is visibly under revision: a blacked-out earlier functor label is replaced by $K_\pi$, a handwritten $K_\pi$ is itself inserted and cancelled before the typed phrase ``defined via,'' and several category-name marks remain doubtful. The checked text does not reconstruct those older labels from the mathematical context. Page~509 explains why $K_\pi$ is conservative using $\pi_i(K_\pi(L.))\simeq H_i(L.)$, while explicitly denying full faithfulness and the assertion that the counit $LH.(K_\pi(L.))\to L.$ is an isomorphism. The typed letter $E$ in the preceding adjunction sentence is retained even though the surrounding notation uses $K_\pi$.

The unit $X\to K_\pi(LH.(X))$ leads to the Hurewicz maps and a proposed homotopy-fiber interpretation of the Whitehead exact sequence (p.~509). The reduced homology in that sequence immediately calls for care with base points. Diagrams(8) and(10) factor the construction through pointed objects, using free pointing $X\amalg e$, the forgetful functor, and reduced free abelianization (pp.~509--510). The quotient description is replaced by an exact triangle, and the parenthetical typed ``one hopes!'' qualifies the continuing claim of adjunction. At the end of p.~510 the author explicitly corrects the earlier passage: the sequence and its fiber should come from the pointed unit(6'), since the unpointed expression lacks the canonical base points needed to define the indicated homotopy groups and fibers. The scan's(6') uses a tilde on both $K$ and $LH$, with no $\pi$ subscript on $K$; the witness regularizes this differently. Section~122 then revisits reduced abelianization itself, seeking a definition that no longer requires the chosen point.

\paragraph{122. Reduced abelianization without a point, and the failed homotopy-linearization functor (pp. 511--517; 134-7/69--75).}
The entry dated 2.10 revises the pointed construction of section~121. Instead of taking the quotient by the chosen base point, Grothendieck defines $\widetilde W(X)$ as the kernel of the augmentation $\underline Z^{(X)}\to\underline Z$ (p.~511). The fiber $W(X)(1)$ is the associated torsor and receives the universal map from $X$ to an abelian torsor. A point splits the extension, recovering the earlier reduced construction and giving the direct-sum description of homology; without a chosen point the natural exact triangle is arranged differently from the triangle on p.~510 (p.~512). The comparison is more informative over an object $Y$, because a section need not exist and the torsor class $c(X/Y)$ can be nontrivial. This supplies the connection with Postnikov invariants and motivates a more general Kan--Dold--Puppe correspondence for coefficients in the relevant slice topos (p.~513). The generalization to local test categories is announced as an intuition to revisit, not established here. The paper's notation is also simplified: the old $K_\pi$ and its pointed variant become $K$ and $\widetilde K$.

A proposed partner $S$ or $\widetilde S$ would realize prescribed reduced homology, in contrast with the homotopy groups prescribed by $\widetilde K$ (pp.~514--515). The coefficient letter in the sphere-like objects $S(H,n)$ is changed in pen from $\pi$ to $H$, and circled marginal $H$ notes emphasize the change. Restriction to connected complexes addresses the immediate obstruction that degree-zero homology of a space is free. The author admits that he has neither constructed the desired functor nor established its existence or uniqueness. In the schematic setting even the relevant low-dimensional sphere objects have not yet been constructed. The possible value of this program depends on the tentative Lie-functor description from section~118, not on an available computation of sphere homotopy groups.

The proposed fourth functor $L\pi.$ is more sharply rejected. It would turn homotopy groups into homology of a chain complex and satisfy $L\pi.(\widetilde K(L.))\simeq L.$. The continuation dated 3.10 states that these properties cannot all hold (p.~515). Applying the functor to the abelianization unit would provide a canonical section of the counit. On $L.=\pi[n]$, this would split the universal-coefficient sequence naturally in the groups involved (p.~516). The corrected indices matter: formula(18) has $\mathrm{Ext}^1$, $H^{n+1}$ and $H_{n+1}$, where the witness gives different indices. The text then specializes to degree one over $\underline F_2$ and the sequence $0\to V'\to\mathrm{Sym}^2(V')\to\bigwedge^2V'\to0$ (p.~517).

An editorial mathematical qualification is needed at the last step. The manuscript claims equivariant nonsplitting for $\dim V\geq2$; the literal bound is retained. If $\mathrm{Aut}(V)$ means the finite group $\mathrm{Aut}_{\underline F_2}(V)$, as the preceding discussion of automorphisms of the abstract abelian group suggests, dimension two is an exception: $q=x^2+xy+y^2$ is fixed by all six invertible linear substitutions over $\underline F_2$, while its polar form maps to the generator of $\bigwedge^2V'$. It therefore supplies an equivariant section for that finite-group action. The qualification is specific to this interpretation: the calculation does not give equivariance for the algebraic group scheme $\mathrm{GL}(V)$ or for all matrices after scalar extension. This small calculation identifies a limitation of the asserted bound on the finite-group reading and does not repair the manuscript's text. The ensuing remark is conditional in a second way: it assumes that an analogous nonsplitting example can be found for $n\geq2$ before drawing the conclusion for simply connected homotopy types. Grothendieck presents that expected example as an alternative route to the failure of the earlier discrete/schematic comparison, avoiding reliance on the reported Breen results. Section~123 opens by acknowledging that, even after these discarded expectations, the basic schematic homotopy operations and the correct weak equivalences are still not in hand.

\paragraph{123. Sphere complexes and scalar operations as tests of the proposed theory (pp. 518--520; 134-7/76--78).}
The entry dated 4.10 takes stock after the preceding counterarguments. Even the equivalence of the proposed criteria for weak equivalence, and the construction of fibers and cofibers, are not yet settled (p.~518). Grothendieck therefore treats the realization of the ordinary two-sphere over $\underline Z$ as a basic test. Suspension should then produce the higher spheres; a larger category admitting the connected complexes needed for $S(1,k)$ might supply an alternative starting point (pp.~518--519). These are requirements for a workable theory, not constructions already obtained.

If schematic sphere objects and the Lie description of their homotopy invariants exist, their homotopy groups should be the homology of canonical complexes $L\pi.(S^n_{\underline Z})$. Maps between spheres would then induce maps of those complexes, compatibly with composition (p.~519). The qualification from section~122 is crucial: the author has just rejected such a functor on arbitrary ordinary homotopy types. He now asks for a special structure on the sphere subcategory and explicitly allows the possibility that this structure is impossible, which would refute the intended schematic theory. The notation is still rough: the defining display for ${}^n\Pi.$ literally writes $L\,.(S^n)$, omitting the $\pi$ present in the surrounding discussion. The checked text retains that omission and the hand-corrected exponent and punctuation in the preceding displays.

The first operation to examine is the action of the multiplicative monoid of integers, obtained by taking maps $S^n\to S^n$ (p.~520). The author relates the resulting actions on finite homotopy groups to the profinite monoid $\widehat{\underline Z}^{\mathrm{mult}}$ and proposes its lift to the derived base extension ${}^n\Pi.\otimes^{L}_{\underline Z}\widehat{\underline Z}$. The derived tensor symbol in the scan is not a direct-sum symbol, and no parentheses surround the tensor expression. For odd-dimensional spheres, restriction to the units is interpreted through the cyclotomic Galois action on the profinite homotopy type of punctured affine space. The pen changes from ``maximal'' to ``largest'' subgroup, adds ``(complex),'' and brackets the exponent $2m-1$; the above-line ``profinite'' remains typed. The point of the passage is the additional structure carried by these operations and a hoped-for lift to complexes, rather than a calculation of the sphere groups. Section~124 next separates the distinctive features of this schematic approach from the two alternatives considered earlier.

\paragraph{124. Four functors and a provisional formal calculus (pp. 521--530; 134-7/79--88).}
Grothendieck contrasts the present unipotent-bundle models with the earlier de Rham and diagonal-category approaches of sections~94 and~109. All seek a homology functor, but the present construction also paraphrases the Eilenberg--Maclane functor $\widetilde K$ compatibly with extension of scalars (p.~521). Simply forgetting a chain complex's $k$-module structure and taking a de Rham complex would not supply that compatibility. Moreover, the present schematic types map to ordinary types, whereas the earlier approaches had the reverse direction of comparison. The notation $\mathrm{Hotab}_0(k)$ deliberately leaves its algebraic realization open: section~115 required complexes of proobjects, and taking their inverse limits would lose the asserted base-change behavior. Consequently $\widetilde K$ may require pro-unipotent models on the nonlinear side as well. Pointed linearization offers a different formal theory, whose invariants section~117 has already distinguished from schematic ones (p.~522).

The continuation dated 5.10 organizes four expected functors between $\mathrm{Hotab}_0(k)$ and $\mathrm{Hot}_0(k)$: total homology $L\widetilde H.$, total homotopy $L\pi.$, and their partners $\widetilde K$ and $\widetilde S$ (p.~523). Even the zero-connected convention is provisional: the author lacks an intrinsic satisfactory formulation of abelian fundamental group with trivial higher action, and briefly considers restricting to simply connected objects. Property~A posits the adjoint pairs $(L\widetilde H.,\widetilde K)$ and $(\widetilde S,L\pi.)$. Their units and counits give linearization and colinearization. The blank after $L$ in formula(5)b remains unfilled in the scan; the checked text does not silently add $\pi$. The Hom formulas interpret $K(\pi,n)$ as representing cohomology and $S(H,n)$ as representing mixed homotopy with a ``co-coefficient'' $H$. Taking $H=k$ recovers the expected ordinary homotopy-module functor (p.~524). This is the stated motivation for the spherical adjunction, whose existence has not been established. The circled marginal $H$ and $\widetilde W(X)$ record changes to the coefficient and linearization notation.

Property~B gives the inversion identities $L\pi.(\widetilde K(L.))\simeq L.$ and $L\widetilde H.(\widetilde S(L.))\simeq L.$ (p.~525). Property~C obtains a Hurewicz map by either applying total homotopy to the linearization map or total homology to the colinearization map. The desired compatibility is that these constructions agree. The first nonvanishing homology objects should occur in the same dimension, with an isomorphism there; a pen addition also demands an epimorphism in the next dimension. The lengthy reduction to Eilenberg--Maclane homology in the preceding paragraph is present in the typescript even though the witness omits it. Property~D states preservation of cofibration sequences by the left adjoints and of fibration sequences by the right adjoints. The two exact triangles spell out the homology and homotopy cases (pp.~525--526). In the second triangle the scan's corrected top vertex is $L\Pi.(Y)$ and its bottom row is $L\Pi.(Z)\to L\Pi.(X)$; the witness gives a different cyclic arrangement. Both oblique labels are the numeral~1.

The exactness discussion reveals an unwanted asymmetry. Suspension on the connected linear category is a shift and is fully faithful, whereas loops require truncation and lose information (p.~526). The first truncation bound has been corrected from~0 to~1. Remark~1 explains that the loop and fiber constructions therefore mean the neutral or marked connected component of the unrestricted construction. Grothendieck says that he has only just become aware of this restriction while writing, questions the practical value of imposing $X_0=e$, and contemplates enlarging both categories (p.~527). The wished-for trivial action of $\pi_1$ on higher groups is a different condition and is not supplied by $X_0=e$. Remark~2 recalls the Leray spectral sequence for homology of fibrations and asks whether a corresponding mechanism exists for homotopy of cofibers; it does not supply such a mechanism.

Remark~3 starts from the evident exchange symmetry of the four-functor formulas, but then distinguishes autodual axioms from an actual autoduality of the categories (pp.~527--529). Full faithfulness of suspension must be excluded from the symmetric axioms. Derived duality on perfect complexes sends chains to cochains, prompting a tentative appeal to stabilization. The subsequent recollection of geometric duality is explicitly qualified by ``as far as I remember'' and ``if I got it right''; the former qualification, the word ``discovered,'' and the shift qualification are typed insertions, placed with pen brackets (p.~529). That discussion pertains to homology and still does not produce a duality exchanging homology and homotopy. The literal sentence ``homotopy type of X'' omits the vee used on the complement elsewhere; the checked text preserves the omission.

Properties~E--G return to the operational requirements (pp.~529--530). Homology and homotopy should be conservative in the indicated connected setting, and all four functors together with their structural maps should commute with extension of the ground ring. Restriction of scalars is mentioned without a developed list of formal properties. Comparison with ordinary types combines forgetting $k$ on the linear side with taking sections of schematic models. It should respect $\widetilde K$, but the already disproved homology comparison is explicitly excluded, even for $k=\underline Z$. Compatibility with $\widetilde S$ is still hoped for over $\underline Z$, conditional on constructing the discrete spherical functor. There is no ordinary total-homotopy functor to compare with $L\pi.$; only the separate homotopy invariants are available there. The closing expectation that the comparison respects fiber and cofiber structures is another requirement of the program. Section~125 begins on13.10 with Baues's proposed axiomatic setting for making those operations precise.

\paragraph{125. From Baues's categories to weak equivalences and derivators (pp. 530--532; 134-7/88--90).}
The entry dated13.10 records a letter from H.~J.~Baues and a reading of his preprint on the homotopy classification problem. Grothendieck recognizes the relevance of the dual notions of cofibration and fibration category to his integration and cointegration operations (p.~530). The present semisimplicial unipotent models, however, do not admit the pushouts he would need for a cofibration structure, except in the stated trivial case. He has begun to consider a Kan-fibration notion but reports no definite conclusion. He therefore proposes to leave the schematic construction temporarily and return after further work on these categorical methods (p.~531).

The more general objection concerns the scope of the axioms. Grothendieck questions whether every meaningful homotopy universe must arise from Quillen's or Baues's extra structures on a chosen category of models. He invokes semisimplicial objects in a topos, Illusie's weak equivalences, and Verdier's derived categories as reasons to seek a wider formulation (p.~531). These are the author's contemporary judgments and questions, not proofs of general nonexistence of model structures. The qualifications ``it seems,'' ``I would think,'' and ``but (it seems)'' are significant: the last is a pen revision to the claim about enough projectives or injectives and bounded derived categories (p.~532). The empty reference brackets after Illusie and the marginal cross remain empty in the text. The circled $C^{-}(A)$ and $D^{-}(A)$ correct the signs used in the discussion of the bounded categories.

The proposed replacement starts with a category $M$ and weak equivalences $W$, without treating a selected fibration or cofibration class as the fundamental datum (p.~532). Different compatible choices should yield the same integration and cointegration operations. Section~69 supplies the earlier point of reference; a derivator is now identified as the structure that should express those operations independently of a particular model presentation. Constructing and characterizing the schematic derivator remain hopes. The passage also frames the derivator as the further structure missing when one stops at a localized or triangulated category, in both abelian and nonabelian settings. The typed insertions ``specific,'' ``on it,'' and ``make any mention of models'' sharpen that distinction, while the correction from ``is'' to the typed ``are'' refers to the two operations. Section~126 resumes the schematic discussion by testing the formal symmetry against Postnikov decomposition.

\paragraph{126. Postnikov truncation and the failure of a literal dual filtration (pp. 533--535; 134-7/91--93).}
The section tests the four-functor symmetry on the construction that originally motivated the schematic program. An $n$-co-connected object has no homotopy groups above degree $n$, and the proposed truncation $X\to X_n$ is the universal map to that full subcategory (p.~533). Equivalently, its inclusion should admit a left adjoint. This universal property produces the familiar inverse tower, with $X$ viewed as an appropriate inverse limit (p.~534). The discussion is axiomatic: it asks what should follow from the basic functors and their structural maps. The double underline on ``co'' distinguishes the upper-degree vanishing condition from connectedness.

A literal dual would be a canonical increasing filtration ${}_nX\to X$, inducing the required homology isomorphisms through degree $n$ and vanishing above that degree. The ordinary skeletal filtration misses the endpoint of the desired isomorphism range. After trying to modify it, Grothendieck reports the feeling that no such canonical filtration exists, while explicitly admitting that he has neither formulated nor proved a nonexistence statement (p.~534). The decisive structural objection is that both proposed truncation subcategories are inverse images of the same linear subcategory under $L\pi.$ or $LH.$. That linear subcategory is stable under loops but not under suspension (p.~535). Even if the doubtful right adjoints existed on the homology side, their successive terms would not have the literally dual Postnikov relationship. A different suspension-stable family is contemplated and then abandoned for lack of a reasonable candidate.

The conclusion qualifies the formal symmetry of section~124: homology and homotopy retain materially asymmetric properties. Several visible details prevent the reader from mistaking the formulas for a fully regularized exposition. On p.~533 the object named $Y$ is tested by a formula using $X$; on p.~534 the skeletal vanishing display writes $H_i(X_n)$ with a right subscript despite the surrounding left-subscript filtration. Formula labels are changed by hand, while some references to(5) remain on p.~535 after the relevant subcategory has been renumbered. Those literal features are retained. Section~127 then changes the question from symmetry of operations to the choice of components allowed in a schematic model.

\paragraph{127. Which objects should serve as schematic models? (pp. 535--538; 134-7/93--96).}
The section keeps the usual simplicial indexing category and changes the category of allowed components. The broad ambient setting is the fpqc topos of sheaves over the ground ring, with Illusie's weak equivalences of semisimplicial sheaves (pp.~535--536). Such a weak equivalence need not induce a weak equivalence on ordinary sections. Grothendieck therefore distinguishes the section functor from the canonical cointegration operation that should recover a discrete homotopy type. The above-line words ``canonical'' and ``(= weak equivalence)'' are typed revisions. He hopes that a suitably restricted class of models, especially the earlier unipotent bundles, will make the naive section construction adequate. The handwritten sentence opening p.~537 makes the base-change requirement explicit: a schematic homotopy type over $k$ should give a functor from $k$-algebras to ordinary homotopy types.

The relative homotopy sheaves are defined first after choosing a section. The proposed ``pseudo-abelian'' condition combines relative zero-connectedness, abelian fundamental group, and trivial action on the higher homotopy sheaves (p.~537). Under a local formulation of that condition, the abelian sheaves should be independent of the chosen section and exist even without a global section. Their desired quasicoherent module structures are an additional issue. The comparison from the homotopy of sections to sections of the homotopy sheaves should be an isomorphism under mild hypotheses, for example $X_0=X_1=e$; the author expressly says he has not attempted a proof for unipotent bundles. He still calls the module structure a mystery without the Lie functor of section~118. The final remark proposes a spectral sequence from cohomology of the homotopy sheaves to the homotopy of cointegration, suggesting that acyclicity might imply the comparison (pp.~537--538). This remains a suggested argument, marked by ``maybe'' and ``If so,'' rather than a constructed spectral sequence. Section~128 next imposes a separate closure requirement: linearization should itself remain in the chosen category of bundles.

\paragraph{128. Linearization and the competing $W$- and $V$-bundles (pp. 538--541; 134-7/96--99).}
The new requirement is that linearization of an allowed bundle should itself be an allowed bundle, carrying the additional structure of an $\underline O_k$-module (p.~538). A forgetful functor $K:Uab(k)\to U(k)$ should therefore have a left adjoint $L$. The two natural choices have opposite variance: $W(M)(k')=M\otimes_k k'$ gives quasicoherent modules, while $V(M)(k')=\operatorname{Hom}_k(M,k')$ gives the opposite of the category of modules. The typed wording does not introduce an extra $X$ after ``every bundle'' or ``for a given bundle,'' despite the witness's additions, and formula(7) has a semicolon after $K$. On passing to semisimplicial objects, these functors are intended to produce the earlier $\widetilde K$ and $LH.$. The final dot of $LH.$ on p.~539 is baseline punctuation in the source, not a supplied subscript.

Section~115 had already shown that the $W$-case gives only a proadjoint. Grothendieck now draws out the cost: one must work with prounipotent bundles and semisimplicial prosheaves outside the sheaf-theoretic foundations he is using (p.~539). A diagonal marginal addition offers working with the inverse limit of the corresponding sheaves; its wording is ``we may however,'' and ``corr.'' remains abbreviated. The witness omits this qualification. The $V$-case is more tractable because $V(M)=\operatorname{Spec}(\operatorname{Sym}_kM)$. Pullback--pushforward adjunction and quasicoherence of the direct image for a quasicompact, quasiseparated scheme give $L_V(X)=V(\Gamma(X,\underline O_X))$, retaining only the underlying module of the global function algebra (pp.~539--540). The Hom subscript in the first display on p.~540 is $\underline O_S$, whereas the witness prints $\underline O_X$. Two circled marginal formulas emphasize the global-sections description. The replacement of ``inclusion'' by the typed ``forgetful'' makes the functor under discussion precise.

This advantage raises a larger question: why restrict components to schemes isomorphic to vector bundles instead of allowing all affine schemes? In the affine setting, semisimplicial schemes correspond to cosemisimplicial commutative algebras, and linearization simply forgets multiplication (p.~540). The answer returns to the original ambition over $\underline Z$, where the schematic category should remain close to ordinary homotopy types. Grothendieck calls the initial choice of $W$ rather than $V$ almost a matter of chance, but formulates two substantive tests for either choice (p.~541). First, the homotopy-sheaf comparison with sections should hold, possibly under $X_0=e$ or $X_0=X_1=e$. Second, under flatness assumptions, the Lie functor should describe the homotopy sheaves. Postnikov d\'evissage supplies only the tentative special case, without the required naturality. He regards recovery from a first-order neighborhood as plausible for affine spaces and implausible for general affine schemes, citing an elliptic curve minus a point. Thus these tests motivate the restricted components; they are not established by the easier linearization adjunction. The next day's section~129 adds the further problem of descending that adjunction to localized categories.

\paragraph{129. The missing descent of linearization and q.c.h. equivalences (pp. 542--543; 134-7/100--101).}
The entry dated14.10 corrects an omission in the preceding day's argument: existence of $L:U(k)\to Uab(k)$ does not itself imply a functor on the corresponding localized categories. One must specify weak equivalences on the linear objects and prove that $L$ preserves them (p.~542). Grothendieck expects preservation and reflection for the $W$- and $V$-models under suitable connectedness assumptions, but explicitly says that a proof has not even been attempted in the unipotent case. This is the third test, after the two tests on p.~541. He compares it with the Whitehead problem already considered in section~92 for replacing the simplex category by other indexing categories. The first live initial before Whitehead in the scan is a corrected typed $J$, over a doubtful older impression. The raised qualifications ``maps between'' and ``total'' are typed additions, and the filled versus open round indices remain graphically distinguished.

For broader classes of semisimplicial schemes, the section proposes defining ``quasi-coherent homological quasi-isomorphisms'' by their image under linearization (p.~543). Ordinary weak equivalence need not yet be known to imply this condition, and a homological criterion can be weaker without connectedness restrictions. The proposed refinement allows twisted quasicoherent coefficients: a sheaf $F_*$ must be carried isomorphically to itself by the structural pullback maps. With a section and relative zero-connectedness, this is related to a fundamental-group-sheaf action on coefficients from the base. The circled marginal abbreviation q.c.h. is present in the manuscript and absent from the witness. Grothendieck then hopes that localization at these maps could make a mild category of general semisimplicial schemes equivalent to the category obtained from bundles, producing a subtler connection with ordinary homotopy types than the known \`etale pro-homotopy connection. He immediately suspends the suggestion as possibly out of reach or mistaken. The typed insertions ``very mild'' and ``only'' emphasize the proposed enlargement and comparison, not a proved equivalence. Section~130 returns to a more limited enlargement by allowing disconnected bundles.

\paragraph{130. Disconnected components as a possible aid to integration (pp. 544--545; 134-7/102--103).}
The limited enlargement now proposed allows direct sums of unipotent bundles as the components of a semisimplicial model. Deleting ``finite'' and replacing the following sentence make the role of infinite sums explicit: linearization still makes sense if the linear category is closed under them; otherwise one must retain a finiteness restriction (p.~544). Descent would also call for twisted finite sums, but the author postpones those technicalities. Applying connected components degree by degree produces a discrete semisimplicial object $E_* = \pi_{o/k}(X_*)$ and a map $X_*\to E_*$. Here $\operatorname{Spec}(k)$ is assumed connected and the sums untwisted. The ``strict'' schematic part lies over this discrete homotopy type, whose own invariants are not themselves schematic. To retain the original strict theory, $E_*$ should be aspheric, explicitly meaning isomorphic to a point in $\mathrm{Hot}$. Equivalence of the resulting localized category with the earlier connected-model category is still only presumed.

The motivation is operational: the ordinary construction of integration from section~69 introduces many connected components, and standard simplices would become available among the models (p.~545). For an aspheric indexing category, Grothendieck hopes to preserve the asphericity of $E_*$. He proposes performing integration and cointegration first among semisimplicial sheaves, then checking whether the results lie in the restricted bundle class. He reports no clear-cut result. Moreover, disconnected components threaten the tangent-sheaf description of section~118, so the enlargement would be justified only after proving that it gives the same localized category as the connected models; no attempt at that proof has yet been made. The typed insertions ``components of the,'' ``sheaf,'' and ``would'' refine this cautious formulation. The cancelled prospective number131 belongs to the continued warning, while the actual section131 starts below a literal date16.9. That date is retained despite the preceding entry dated14.10. Its opening personal interlude records the kimchi instructions and a letter to Larry Breen before the next page sharpens the mathematical objection to the enlargement.

\paragraph{131. Infinitesimal spheres and the failure of the naive Lie formula (pp. 545--551; 134-7/103--109).}
After the personal interlude opening the section, Grothendieck rejects the disconnected-model proposal as largely a symptom of his inability to construct the required ``left'' homotopy operations within connected unipotent bundles (pp.~545--546). The implicit condition $X_o=e$, or even $X_o=X_1=e$, is essential to the expected Lie description. Without it, tangent invariants at different sections need not be canonically identified, already for a constant semisimplicial bundle. Reimposing it would undermine the ordinary integration construction that motivated disconnected components in the first place. The author therefore distinguishes the restricted models for schematic homotopy types from more general semisimplicial sheaves that are useful auxiliary objects. The added pen $o$ in ``loosing track'' is clear and is retained, although the witness prints ``losing.''

A different enlargement is suggested by the still hypothetical spherical functor $\widetilde S$ of section~124, which should be left adjoint to total homotopy and right inverse to reduced total homology (pp.~547--548). For flat augmented symmetric algebras representing pointed $V$-bundles, the augmentation ideal modulo its square gives the co-Lie module. A module $M$ also defines the square-zero algebra $D(M)=k\oplus M$ and its spectrum $I(M)$. Applied degreewise, this construction gives the formal candidate $\widetilde S(L_*)=I(L^*)$ for $L_*=V(L^*)$. The resulting Hom identity expresses ordinary representability of the tangent functor by first-order infinitesimal schemes. Grothendieck explicitly marks two gaps: it is an identity between model objects, not yet an adjunction of derived categories, and the resulting infinitesimal schemes are outside the existing bundle category. The pen correction of ``adjunction'' and the deleted claim that the first Hom is taken in $\mathrm{Hotab}$ make the first gap especially visible on p.~548. Page~547 also literally introduces ``L'' as the arbitrary module immediately before writing $D(M)$; that variable mismatch is retained.

The inversion identity with reduced homology is equally tautological for square-zero algebras, motivating a new definition of bundles as nilpotent thickenings of $V$-bundles (p.~549). The chosen nilpotent ideal is not part of the structure, and further flatness conditions remain possible. The subscheme in the displayed inclusion is $X_o$, whereas the witness prints $X_*$. Three of the functors now have simple model descriptions: augmentation ideals for reduced homology, symmetric algebras for $\widetilde K$, and their first-order truncations for $\widetilde S$. The canonical closed immersion $\widetilde S(L_*)\to\widetilde K(L_*)$ reveals a missing compatibility axiom in section~124 (p.~550). Its two adjuncts should be the inversion isomorphisms, and the maps obtained from those two isomorphisms must agree. Taking this natural transformation as primary data would also give a single Hurewicz map. The literal condition near the foot of p.~550 has right-hand side $k$, with a doubtful round left index; the witness replaces it by $X_0=e$. No such mathematical correction is made in the text.

The apparent simplicity breaks down at the interpretation of total homotopy (p.~551). On the algebra side, naive $L\Pi.$ takes an augmentation ideal to its quotient by the square; dualizing gives the inclusion form of the Hurewicz map. But for a first-order thickening the square is already zero, so $L\Pi.(\widetilde S(L_*))\simeq L_*$. If these were the desired spheres, their naive Lie homotopy would therefore equal their homology, which the author emphatically rejects. His remaining hope is restricted to actual flat vector-bundle components, the basepoint condition, and possibly unspecified Kan hypotheses. The triple exclamation and the two marginal corner marks emphasize the failure and the beginning of section~132. That next section replaces the naive tangent functor by the Andr\'e--Quillen--Illusie tangent and cotangent complexes.
